\documentclass[10pt,twocolumn]{article}

% --- Packages (all standard; compiles with a basic TeX Live / pdflatex) ---
\usepackage[utf8]{inputenc}
\usepackage[T1]{fontenc}
\usepackage[hmargin=1.8cm,vmargin=2.0cm]{geometry}
\usepackage{times}
\usepackage{url}
\usepackage{hyperref}
\hypersetup{
  colorlinks=true,
  linkcolor=black,
  citecolor=black,
  urlcolor=blue,
  pdftitle={arxivMedia: A Social Network Where AI Agents and Humans Review the arXiv Firehose},
  pdfauthor={Jay Desai}
}
\usepackage{booktabs}
\usepackage{enumitem}
\usepackage[numbers]{natbib}
\usepackage{balance} % optional; balances columns on the last page

\setlength{\parskip}{0.3em}
\setlength{\columnsep}{0.7cm}

% --- Title ---
\title{\textbf{arxivMedia}: A Social Network Where AI Agents and\\
Humans Review the arXiv Firehose}

% Sole author. Affiliation is a placeholder for the author to fill before
% any submission; ORCID unknown.
\author{
  Jay Desai\thanks{Affiliation: \texttt{<placeholder --- author to fill before submission>}.}\\
  \texttt{jmd.desai08@gmail.com}
}
\date{June 2026}

\begin{document}
\maketitle

\begin{abstract}
The volume of papers posted to arXiv in machine learning and adjacent fields
now far exceeds what any individual---or any editorial team---can read, triage,
or review. We present \textbf{arxivMedia}, an open-source social network that
treats this firehose as a feed to be reviewed continuously, and that makes
\emph{autonomous AI agents first-class participants} alongside humans. A system
crawler ingests new arXiv submissions and posts one thread per paper; AI agents
register through an open write API and humans sign up through a web interface,
and both write reviews, debate methodology, reply in threads, and vote on the
same engine. Agents onboard themselves by fetching a single machine-readable
document at \texttt{/skill.md} that documents the entire API. The system ranks
papers with Hacker-News-style hot ranking plus new, trending (comment-velocity),
time-windowed top, and most-cited orderings, and enriches papers with external
citation counts from Semantic Scholar. arxivMedia is built on a deliberately
small stack (FastAPI, stdlib \texttt{sqlite3}, Jinja2, no JavaScript build) and
runs on free-tier hosting; a live deployment is available at
\url{https://djaym7-arxivmedia.hf.space}, with reference reviewer agents for
Anthropic Claude and Google Gemini. This is a demonstration system: we describe
its design and position it against existing platforms, but report no user study
or quantitative evaluation. Its novelty is a \emph{configuration}---agents-first
participation, an open self-serve write API, social voting, firehose-scale
triage, and open source---rather than any single new capability. We discuss the
honest open question it raises: whether machine peer review, conducted in the
open, surfaces real signal.
\end{abstract}

% ====================================================================
\section{Introduction}
% ====================================================================

Thousands of papers are submitted to arXiv every day, and a large and growing
fraction of them are in machine learning and artificial intelligence. The rate
of submission in these areas has long since outpaced the capacity of any human
to read even the abstracts, let alone evaluate the claims. Curated venues help:
Hugging Face's Daily Papers surfaces a small daily selection, and conference
peer review evaluates a filtered subset months after preprints appear. But the
\emph{firehose}---the continuous stream of everything submitted in a category---
goes largely un-triaged in public. No human reads it all; no human even tries.

At the same time, language-model agents have become capable enough to read an
abstract, identify a concrete strength and a concrete concern, connect a result
to prior work, flag an overclaim or a likely contamination issue, and articulate
all of this in a few sentences. An agent can do this for \emph{every} abstract in
a category, around the clock, at a cost that free-tier model access already
covers. The natural question is whether a community of such agents, reviewing in
public and arguing with one another, can surface useful signal from the firehose
that would otherwise be lost.

\textbf{arxivMedia} is a small, concrete bet that the answer is worth testing. It
is an open-source social network---structured like Hacker News or lobste.rs---in
which the primary participants are autonomous agents, with humans participating
on the same footing. A system crawler ingests new arXiv papers and posts one
discussion thread per paper. Agents register through an open API, fetch a
machine-readable onboarding document, and then post, review (as threaded
comments), reply, and vote. Humans sign up through the web UI and do exactly the
same things through their browser. The design is directly inspired by
Moltbook~\cite{moltbook}, the ``agents post, humans watch'' social network, but
is scoped specifically to peer review of machine-science papers and is open
source and self-hostable.

This paper is an honest systems/demonstration report. The system exists, is
deployed, and works; but we have run \emph{no} user study and report \emph{no}
quantitative evaluation of review quality. We describe what is built and position
it against related work, and we are explicit about what remains an open question.

\paragraph{Contributions.} This paper contributes:
\begin{itemize}[leftmargin=1.2em,itemsep=0.15em,topsep=0.2em]
  \item \textbf{The system.} A working, open-source, HN-style social network for
  reviewing the arXiv firehose, with a unified account model and shared
  post/comment/vote/ranking engine for agents and humans
  (Section~\ref{sec:design}).
  \item \textbf{An open agent write API with self-serve onboarding.} Any agent,
  anywhere, can register in one request and learn the entire API by fetching a
  single machine-readable document at \texttt{/skill.md}
  (Section~\ref{sec:design}, Section~\ref{sec:impl}).
  \item \textbf{Agents as first-class participants.} A design in which autonomous
  agents are the primary social actors of a public, votable feed---not assistants
  bolted onto a human-driven site---with humans joining on equal terms.
  \item \textbf{A live, free-tier deployment.} A running instance at
  \url{https://djaym7-arxivmedia.hf.space} on Hugging Face Spaces, with reference
  reviewer agents for Anthropic Claude and Google Gemini (Section~\ref{sec:impl}).
\end{itemize}

% ====================================================================
\section{Related Work}
\label{sec:related}
% ====================================================================

arxivMedia sits at the intersection of paper-discovery platforms, peer-review
systems, citation tooling, and the emerging space of agent-native social
networks. We survey each and are explicit about where arxivMedia is, and is not,
novel. The comparison below is drawn from a competitive analysis maintained
alongside the codebase.

\paragraph{Curated paper-discovery platforms.}
\textbf{Hugging Face Papers}~\cite{hfpapers} (including Daily Papers) is the
closest paper-specific platform: it already ingests arXiv submissions, supports
upvotes, author-claiming, and threaded comments, and even runs an assistive bot.
But its social layer is \emph{human-driven}, and it exposes only a \emph{read}
API (e.g.\ \texttt{/papers/\{id\}.md})---there is no public \emph{write} API for
an agent to post, comment, or vote. \textbf{alphaXiv}~\cite{alphaxiv} adds an
open, human-driven social commentary layer over arXiv preprints, with an AI
assistant for reading rather than autonomous posting.
\textbf{Papers with Code} curated benchmark leaderboards linking tasks, datasets,
metrics, and code, but was retired in 2025. \textbf{arxiv-sanity}~\cite{arxivsanity}
provides personalized arXiv recommendation and is open source, but offers
personalized ranking and saved libraries rather than a social feed with votes and
threads.

\paragraph{Peer review and AI reviewing.}
\textbf{OpenReview}~\cite{openreview} manages gated conference and journal peer
review and, importantly, exposes a genuine \emph{read+write} API that AI-review
agents already use~\cite{openreviewer}. But it is gated peer review tied to
specific venues, not an open, public, votable feed, and it is not ``agents
first.'' A separate line of work \emph{generates} reviews or whole papers with
LLMs---OpenReviewer~\cite{openreviewer} and Sakana's AI Scientist~\cite{aiscientist}
among them---but these are pipelines or models that produce one-shot review
artifacts, not a continuous, public, participatory review network.

\paragraph{Citation and reading tools.}
\textbf{Semantic Scholar}~\cite{semanticscholar} offers scholarly search, an
AI-augmented reader, and a free Academic Graph API; arxivMedia \emph{consumes}
its citation-count endpoint to enrich papers (Section~\ref{sec:design}), but
Semantic Scholar itself has no social posting layer. Citation-graph tools such as
Connected Papers and citation-context indexes such as Scite.ai augment reading
and literature mapping but, again, are not participatory social feeds.

\paragraph{Agent-native social networks.}
\textbf{Moltbook}~\cite{moltbook} is the direct conceptual parent: a social
network where \emph{agents} post, comment, and upvote while humans only watch.
arxivMedia adopts that ``agents-first'' inversion but scopes it to arXiv paper
review, adds an open write API and machine-readable onboarding, lets humans
participate (not just watch), and is open source---whereas Moltbook is a closed
product and is about open-ended discussion rather than papers.

\paragraph{Where arxivMedia is, and is not, novel.}
Every \emph{individual} capability here exists somewhere already: arXiv ingestion
(Hugging Face Papers), social voting and comments on papers (Hugging Face Papers,
alphaXiv), AI-generated reviews (OpenReviewer, AI Scientist), a read+write review
API (OpenReview), citation enrichment (Semantic Scholar), and ``agents post,
humans watch'' (Moltbook). arxivMedia is novel as a \emph{configuration}: a
public, open-source, votable paper feed in which autonomous agents are the
primary participants via an \emph{open self-serve write API} with machine-readable
onboarding, doing continuous triage and review of the \emph{full arXiv firehose},
with humans participating on the same engine. No existing platform combines all
of these. We state this plainly: the contribution is a product configuration and
a working demonstration, not a brand-new capability.

% ====================================================================
\section{System Design}
\label{sec:design}
% ====================================================================

arxivMedia is a single server-rendered web application that also exposes a JSON
API. Its components are intentionally few.

\paragraph{Ingestion and deduplication.}
A system account, \texttt{arxiv-crawler}, periodically queries the arXiv Atom API
for the most recent submissions in a configurable set of categories (by default
\texttt{cs.CL}, \texttt{cs.LG}, \texttt{cs.AI}). The Atom XML is parsed with the
standard library, titles and abstracts are whitespace-normalized, and each paper
becomes one post whose \texttt{source} field is set to \texttt{arxiv:\{id\}} (the
version suffix stripped). A uniqueness constraint on \texttt{source} deduplicates
papers by arXiv ID, so re-running ingestion never creates duplicate threads.

\paragraph{Unified accounts.}
Agents, humans, and the system crawler are all rows in a single \texttt{agents}
table distinguished by a \texttt{kind} field (\texttt{agent}, \texttt{human}, or
\texttt{system}). Agents authenticate with an API key sent in an
\texttt{X-API-Key} header; humans authenticate with a password (PBKDF2-HMAC-SHA256,
salted, stored as a hash) and a signed session cookie. Crucially, the
create-post, create-comment, and cast-vote logic is shared: both the JSON API and
the human web routes call the same helper functions, so there is exactly one
implementation of the social mechanics and one karma/scoring path for everyone.

\paragraph{Posts, threaded comments, and votes.}
The data model has posts (paper threads or agent-authored discussions), comments
(reviews and replies, nested via a \texttt{parent\_id} self-reference into a
tree), and votes (an upsert keyed by account, target type, and target id; re-casting
the same vote toggles it off). A vote applies a score delta to the target and a
karma delta to the target's author; accounts cannot vote on their own content.
Per-process, in-memory rate limits cap requests, registrations, posts, and
comments.

\paragraph{Ranking and discovery.}
The feed supports five orderings: \emph{hot} (a Hacker-News-style time-decayed
score, computed in Python over recent posts), \emph{new} (reverse-chronological),
\emph{trending} (recent comment-velocity plus a fraction of total score over a
fixed 48-hour engagement window, decaying more slowly than hot), \emph{top}
(highest score within a \texttt{day}/\texttt{week}/\texttt{month}/\texttt{all}
window), and \emph{cited} (most external citations first). A single \texttt{area}
parameter filters any ordering to one arXiv category, and a \texttt{GET
/api/areas} endpoint lists categories with post counts so clients can build a
selector. Invalid sort/window values are clamped to defaults rather than
returning errors.

\paragraph{Citation enrichment.}
Papers are enriched with external citation counts from the Semantic Scholar Graph
API, looked up by arXiv ID. All such network I/O happens in the background
ingestion thread---never on the request path---paced to roughly one request per
second and capped per cycle, with results cached in two nullable columns on the
post row and periodically refreshed. Every failure (404, rate limit, timeout)
degrades gracefully to ``no data'' and never aborts ingestion.

\paragraph{Machine-readable onboarding.}
\texttt{GET /skill.md} returns a single Markdown document, with the live base URL
substituted in, that documents every endpoint with copy-pasteable
examples---registration, authentication, reading the feed, posting, commenting,
voting, rate limits, and reviewing etiquette. An agent can fetch this one document
and have everything it needs to participate, with no human in the loop.

\paragraph{Web interface.}
The human-facing UI is a dark, dense, HN/lobste.rs-style interface rendered
entirely server-side with Jinja2 templates and a small static stylesheet---no
JavaScript build step, no client framework. Humans browse the ranked feed, read
paper threads with nested reviews, and (when logged in) submit, comment, reply,
and vote through plain HTML forms.

\paragraph{Stack and persistence.}
The system is built on FastAPI with the standard-library \texttt{sqlite3} module
(no ORM) and Jinja2. For deployment on an ephemeral filesystem, an optional
persistence layer snapshots the SQLite database to a Hugging Face Dataset on a
timer (using SQLite's online-backup API for a consistent copy) and restores it on
boot, so live data survives restarts. This component is a no-op in local
development and never raises into the request path.

% Figure placeholder (optional; kept simple so the paper compiles without an
% external image). Uncomment and supply a file to include a real figure.
% \begin{figure}[t]
%   \centering
%   \includegraphics[width=\columnwidth]{architecture.png}
%   \caption{arxivMedia architecture: arXiv ingestion feeds a shared
%   post/comment/vote engine consumed by agents (JSON API) and humans (web UI).}
%   \label{fig:arch}
% \end{figure}

% ====================================================================
\section{Implementation \& Deployment}
\label{sec:impl}
% ====================================================================

The implementation is deliberately compact: roughly 1{,}400 lines of Python
across a handful of modules (the FastAPI app and routes, the SQLite schema and
migrations, arXiv ingestion, citation lookups, and the optional HF-Dataset
persistence), plus a small set of Jinja2 templates and one stylesheet. The only
runtime dependencies are FastAPI/uvicorn, Jinja2, \texttt{httpx} (for ingestion
and citation lookups), and the session/form-parsing helpers; everything else is
the Python standard library. There is no JavaScript build, no ORM, and no
external font or CDN dependency.

\paragraph{Deployment.}
A live instance runs on \textbf{Hugging Face Spaces} (free tier) at
\url{https://djaym7-arxivmedia.hf.space}, packaged as a small Docker image. The
ingestion loop begins on startup, so a freshly deployed instance populates its
feed within minutes. The repository also ships configuration for Render and
Fly.io and a generic Dockerfile for self-hosting. All configuration is via
optional environment variables (database path, ingestion categories and interval,
persistence and citation toggles), with sensible defaults.

\paragraph{Reference reviewer agents.}
Two reference agents are provided, both powered by free-tier model access. One
uses Anthropic's Claude and one uses Google's Gemini Flash. Each registers (or
reuses a saved identity) through the API, reads the feed, writes a short scholarly
review of each unseen abstract---prompted to name one concrete strength and one
concern grounded in specifics, in plain text---posts it as a comment, and upvotes
papers it found valuable. These agents demonstrate the intended loop and serve as
templates anyone can adapt; the \texttt{/skill.md} document is sufficient to write
a new one against any model provider.

% ====================================================================
\section{Discussion \& Limitations}
\label{sec:limitations}
% ====================================================================

We are deliberately conservative about what this work shows.

\paragraph{No evaluation yet.}
This is a demonstration system. We have not run a user study, and we report no
quantitative measurement of whether agent reviews are accurate, useful, or better
than no review. Whether machine peer review at firehose scale surfaces real
signal---rather than plausible-sounding noise---is an open empirical question that
this paper raises but does not answer.

\paragraph{Review quality and hallucination.}
Agent reviews are only as good as the underlying models and prompts. An agent can
confidently misstate a paper's method, invent a weakness, or miss a real one.
Because reviews are written from abstracts (and optionally the paper), they are
especially exposed to overclaiming. The system encourages substantive, specific
reviews through its onboarding etiquette, but it does not---and cannot---guarantee
correctness.

\paragraph{Spam and Sybil risk.}
Registration is open and key issuance is unauthenticated by design, which is what
makes self-serve onboarding possible---but it also means an adversary could
register many agents to manufacture votes or flood threads. Current defenses are
limited to per-process rate limits; there is no proof-of-work, identity
verification, or robust karma gating yet. At scale, this is the most pressing
weakness.

\paragraph{Operational and coverage limits.}
The default deployment uses SQLite with snapshot-based persistence to a Hugging
Face Dataset; this is sufficient for a free-tier demonstration but is not a
high-durability production datastore. Deduplication is by exact arXiv ID, so the
same work appearing under different identifiers (e.g.\ a later venue version) is
not merged. Citation coverage depends on Semantic Scholar having indexed a paper,
which newly posted preprints often are not, and on staying within shared-pool
rate limits.

% ====================================================================
\section{Future Work}
\label{sec:future}
% ====================================================================

Several directions follow directly from the limitations above.
\begin{itemize}[leftmargin=1.2em,itemsep=0.15em,topsep=0.2em]
  \item \textbf{Evaluation study.} Compare the review quality of different models
  on an \emph{identical} set of papers---ideally papers with known outcomes or
  expert annotations---to measure accuracy, calibration, and usefulness rather
  than assuming them.
  \item \textbf{Anti-spam and agent verification.} Proof-of-work or staked
  registration, karma gates, and reputation that make Sybil attacks costly, so the
  feed remains a credible signal rather than a bot firehose.
  \item \textbf{Semantic deduplication.} Merge the same work across identifiers
  and sources using embedding similarity rather than exact arXiv-ID matching.
  \item \textbf{Broader ingestion.} Add OpenReview and ACL Anthology streams
  alongside arXiv so the firehose is not arXiv-only.
  \item \textbf{In-network citation graph.} Let agents make structured,
  contestable claims (``this paper claims X'') that other agents support or
  dispute with evidence, building a citation/claim graph internal to the network.
  \item \textbf{Federation.} Allow independent arxivMedia instances---for example,
  a lab running a private instance over its own paper stream---to federate.
\end{itemize}

% ====================================================================
\section{Conclusion}
\label{sec:conclusion}
% ====================================================================

The arXiv firehose has outgrown human triage, and capable agents now exist that
could, in principle, read and review it continuously in the open. arxivMedia is a
working, open-source demonstration of one way to let them: a social network where
autonomous agents are first-class participants via an open self-serve write API,
where humans join on the same engine, and where papers are ranked, reviewed, and
argued over in public. Its novelty is a configuration rather than a new
capability, and we make no evaluative claims about review quality---that is the
central open question we hope the system makes it possible to study. The code is
open source and a live instance is running; we offer it as a substrate for
asking whether machine peer review, in the open, is more useful than no review at
all.

% --- Bibliography ---
\balance
\bibliographystyle{plainnat}
\bibliography{refs}

\end{document}
